% !TeX program = lualatex
% !TeX encoding = utf8
% !TeX spellcheck = uk_UA
% !TeX root = ../main.tex

%=========================================================
\chapter{Явище електромагнітної індукції}\label{\currfilebase}\hypertarget{ElMagInduction}{}
\graphicspath{{\currfilebase/Pictures}}
%=========================================================
\epigraph{\Alexander Магнітне поле не діє на нерухомі заряди~--- і~все ж змінне магнітне поле
	народжує електричне поле, здатне їх розігнати. У~цьому факті криється, мабуть, найглибша
	ідея класичної електродинаміки: поле не просто \enquote{обслуговує} заряди, воно живе
	власним життям і~саме породжує інші поля. }{Переосмислення ідей сучасної фізики}



%% --------------------------------------------------------
\section{Явище електромагнітної індукції}\hypertarget{induction-phenomenon}{}
%% --------------------------------------------------------

%Восени 1831~р. М.~Фарадей здійснив те, що ще в 1822~р. записав собі в пам'ятній книжці як мету --- \enquoteперетворити магнетизм на електрику.
У~1831~р. М.~Фарадей відкрив \emphx{явище електромагнітної індукції}: воно полягає в~тому,
що в~замкненому провідному контурі при зміні магнітного потоку, охопленого цим контуром,
виникає електричний струм, який називають \emphx{індукційним}.

%\begin{figure}[h!]\centering
%	\includegraphics[width=1\linewidth]{Faraday_in_lab}
%	\caption{Майкл Фарадей у своїй лабораторії\label{pic:FaradayLab}}
%\end{figure}

У~низці дослідів Фарадей встановив, що індукційний струм виникає щоразу, коли змінюється
магнітний потік крізь контур~--- незалежно від \emphx{причини} цієї зміни: чи то рухається
магніт відносно нерухомого контуру, чи рухається сам контур відносно нерухомого магніту,
чи змінюється сила струму в~сусідньому контурі (рис.~\ref{pic:FaradayExperiments}).

Розглянемо докладніше ці три групи дослідів.

%% --------------------------------------------------------
\paragraph{Дослід 1. Рух магніту відносно нерухомої котушки.}
%% --------------------------------------------------------

Котушку, з'єднану з~гальванометром, закріплюють нерухомо. До неї наближають або віддаляють
постійний магніт. Поки магніт нерухомий, стрілка гальванометра стоїть на нулі. Щойно
магніт починає рухатися, стрілка відхиляється~--- у~котушці виникає індукційний струм. Під
час наближення магніту струм має один напрям, під час віддалення~--- протилежний. Коли рух
припиняється, струм зникає.

%% --------------------------------------------------------
\paragraph{Дослід 2. Рух котушки відносно нерухомого магніту.}
%% --------------------------------------------------------

Магніт залишають нерухомим, а~рухають котушку, з'єднану з~гальванометром. Спостерігається
той самий ефект: під час руху виникає індукційний струм, напрям якого залежить від напряму
руху котушки. Отже, важливою є не абсолютна швидкість магніту чи котушки, а~їхня
\emphx{відносна} швидкість: саме від неї залежить швидкість зміни магнітного потоку крізь
контур.

\begin{figure}[h!]\centering
	\includegraphics[width=1\linewidth]{Faraday_experiments}
	\caption{Схема класичних дослідів Фарадея з~електромагнітної індукції: виникнення
		індукційного струму при русі магніту відносно котушки (верхній лівий та верхній
		правий), при взаємному переміщенні двох котушок (нижній лівий) та при русі провідника
		в~магнітному полі (нижній правий)\label{pic:FaradayExperiments}}
\end{figure}

%% --------------------------------------------------------
\paragraph{Дослід 3. Зміна струму в~сусідній котушці.}
%% --------------------------------------------------------

Дві котушки розміщують поряд: одну (1) приєднують до гальванометра, другу (2)~--- через
ключ до джерела струму. Коли ключ замикають або розмикають, у~першій котушці виникає
короткочасний індукційний струм. Струм спостерігається лише в~моменти \emphx{зміни} струму
в~другій котушці; коли струм сталий, гальванометр нічого не показує. Це явище називають
\emphx{взаємоіндукцією}.

Узагальнюючи результати всіх цих дослідів, Фарадей дійшов висновку (у~сучасному
формулюванні): \emphx{індукційний струм у~замкненому контурі виникає тоді й~лише тоді,
	коли змінюється магнітний потік, охоплений цим контуром}. Сама ж причина зміни магнітного
потоку~--- рух магніту, рух контуру чи зміна струму в~сусідньому контурі~--- на існування
та напрям індукційного струму принципово не впливає.


%% --------------------------------------------------------
\section{Закон електромагнітної індукції}\hypertarget{induction-law}{}
%% --------------------------------------------------------

На основі дослідів Фарадея було сформульовано \emphx{закон електромагнітної індукції}:
електрорушійна сила (ЕРС), що виникає в~контурі, пропорційна швидкості зміни магнітного
потоку, який пронизує площу, охоплену цим контуром:
\begin{equation}\label{eq:FaradayLawInt}
	\tcbhighmath{ \emf_\text{інд} = -\frac1c\frac{d\Phi}{dt} =
		-\frac1c\frac{d}{dt}\iint\limits_S \Bfield\cdot d\vect S. }
\end{equation}
Знак \enquote{мінус} у~цьому законі~--- не формальність, а~прояв \emphx{правила Ленца}, з
яким ми детальніше познайомимось нижче.

Магнітний потік $\Phi$, а~отже, і~ЕРС індукції, можуть змінюватися трьома способами:
\begin{itemize}
	\item через зміну \emphx{площі} контуру, що пронизується незмінним магнітним полем;
	\item через зміну \emphx{орієнтації} контуру відносно поля (наприклад, при обертанні
	      рамки);
	\item через зміну самого \emphx{магнітного поля}, що пронизує нерухомий контур.
\end{itemize}
У~перших двох випадках контур рухається, у~третьому~--- нерухомий; ці випадки мають різну
фізичну природу, і~ми розглянемо їх окремо.

%% --------------------------------------------------------
\section{Робота, що виконується при переміщенні провідника зі струмом у~магнітному полі}\label{sec:Work_current}\hypertarget{work-of-moving-conductor}{}
%% --------------------------------------------------------

Візьмемо рамку, по якій може вільно ковзати перемичка, що замикає верхній і~нижній проводи
(рис.~\ref{tikz:frame_with_rod}). Помістимо цю систему в~магнітне поле. Будемо вважати
спочатку, що поле перпендикулярне до площини рамки.

\begin{SCfigure}[1.25][h!]
	\localinput{frame_with_rod.tikz}
	\caption{Рамка з~рухомою перемичкою в~магнітному полі: показано силу Ампера
		$\vect{F}_\text{А}$, що діє на перемичку зі струмом $I$, зміщення $dx$ та розклад повної
		швидкості $\vect{u}$ носія заряду на складові $\vect{u}_{\perp}$ (рух разом із перемичкою)
		і~$\vect{u}_{\parallel}$ (рух уздовж провідника)\label{tikz:frame_with_rod}}
\end{SCfigure}



Якщо по рамці через перемичку тече струм \(I\) (як показано на
рис.~\ref{tikz:frame_with_rod}), то на перемичку діє сила Ампера:
\begin{equation}
	F_\text{А} = \frac{I}{c} \ell B, \label{eq:AmpereForceRod}
\end{equation}
де \(\ell\)~--- довжина перемички. При зазначених на рис.~\ref{tikz:frame_with_rod}
напрямках струму \(I\) і~зовнішнього магнітного поля \(\Bfield\) сила
\(\vect{F}_\text{А}\) напрямлена вправо і~при зміщенні перемички на \(dx\) виконує роботу
\begin{equation}
	dA = F_\text{А} dx = \frac{I}{c} B \ell dx = \frac{I}{c} d\Phi. \label{eq:work_dx}
\end{equation}
Площа контуру зі струмом змінюється на \(dS = \ell\,dx\), а~магнітний потік~--- на \(d\Phi
= B\,dS\). При зміщенні перемички на кінцеву відстань (за незмінних \(I\) і~\(\Bfield\))
\begin{equation*}
	A_{12} = \frac{I}{c} (\Phi_2 - \Phi_1).
\end{equation*}

Нехай поле напрямлене під довільним кутом до площини контуру (рамки). Покладемо \(\Bfield
= \Bfield_{\perp} + \Bfield_{\parallel}\), де \(\Bfield_{\perp}\)~--- складова, напрямлена
за нормаллю до площі, а~\(\Bfield_{\parallel}\)~--- що лежить у~площині контуру. Складова
\(\Bfield_{\parallel}\) створює силу, перпендикулярну до площини і, отже, перпендикулярну
до зміщення перемички, і~тому таку, що не виконує роботи. З~урахуванням цього магнітний
потік, що призводить до виконання роботи, визначається тільки компонентою поля
\(\Bfield_{\perp}\):
\begin{equation*}
	\Phi = B_{\perp} S = \Bfield\cdot\vect{S}.
\end{equation*}

У~випадку довільно деформованого витка і~поля, що змінюється від точки до точки, весь
замкнений контур потрібно подумки розбити на нескінченно малі замкнені елементи струму.
Оскільки струми на спільних ділянках сусідніх витків рівні за величиною і~протилежні за
напрямком, то струми по всій площі вихідного контуру компенсують один одного, і
залишається тільки струм \(I\) у~вихідному контурі.


Вважаючи поле в~межах одного елемента однорідним, отримуємо
\begin{equation*}
	dA = \frac{I}{c} d\Phi,
\end{equation*}
де магнітний потік тепер визначається формулою
\begin{equation}
	\Phi = \int\limits_{S} \Bfield\cdot d\vect{S}, \label{eq:magnetic_flux}
\end{equation}
в~якій інтегрування ведеться по площі, що обмежується контуром зі струмом.

Зауважимо, що в~цьому процесі виконується робота, тоді як саме магнітне поле роботу не
може виконувати (сила Лоренца завжди перпендикулярна швидкості руху зарядів). Для
розв'язання цієї суперечності врахуємо, що заряди в~перемичці рухаються не тільки вздовж
провідника (\(\vect{u}_{\parallel}\)), а~й~разом із перемичкою (\(\vect{u}_{\perp}\)):
\(\vect{u} = \vect{u}_{\perp} + \vect{u}_{\parallel}\). Складова \(\vect{u}_{\parallel}\)
породжує силу Лоренца, напрямлену вздовж руху перемички,~--- це і~є сила Ампера
\(\vect{F}_\text{А}\), що виконує додатну роботу над перемичкою. Складова
\(\vect{u}_{\perp}\) породжує силу \(\vect{F}_\text{Л} = q \vect{u}_{\perp} \times \Bfield
/ c\), що діє проти напрямку струму \(I\) і~виконує рівну за модулем від'ємну роботу, так
що сумарна робота сили Лоренца дорівнює нулю. Тому для підтримання незмінного струму в
коло мають бути включені сторонні ЕРС: саме завдяки зовнішнім джерелам енергії виконується
робота з~переміщення перемички.

%% --------------------------------------------------------
\section{Електромагнітна індукція в~рухомих провідниках}\hypertarget{induction-in-moving-conductors}{}
%% --------------------------------------------------------

%% --------------------------------------------------------
\subsection*{ЕРС індукції}
%% --------------------------------------------------------

\begin{SCfigure}[1.25][h!]
	\centering
	\localinput{move_rod.tikz}
	\caption{Виникнення ЕРС індукції в~рухомій перемичці: на позитивний заряд $+$ у
		провіднику, що рухається зі швидкістю $\vect{v}$, діє сила Лоренца $\vect{F}_\text{Л}$,
		спрямована вздовж провідника (швидкість $\vect{u}$). Ця сила викликає індукційний струм
		$I$ у~контурі\label{tikz:move_rod}}
\end{SCfigure}

У~попередньому параграфі ми розглянули роботу сили Ампера й~з'ясували, що рух перемички зі
струмом підтримується сторонніми джерелами. Знайдемо тепер саму ЕРС індукції
\emphx{безпосередньо}, обчисливши силу, що діє на носії заряду: це та сама сила Лоренца
$\vect F_\text{Л}=q\vect u_\perp\times\Bfield/c$, на яку ми вже звертали увагу вище.

Розглянемо провідну рамку, замкнену рухомою перемичкою (див. рис.~\ref{tikz:move_rod}).
Нехай рамка поміщена в~магнітне поле з~індукцією \(\Bfield\), напрямленою перпендикулярно
до площини. Якщо перемичка рухається зі швидкістю \(\vect{v}\) (у~попередніх позначеннях
\(\vect v=\vect u_{\perp}\)), то на заряди в~перемичці діє сила Лоренца:
\begin{equation*}
	\vect{F}_\text{Л} = \frac{q}{c} \vect{v} \times \Bfield.
\end{equation*}
Як видно з~рис.~\ref{tikz:move_rod}, ця сила приводить заряди в~рух, викликаючи
індукційний струм у~від'ємному напрямку обходу контуру. Вона еквівалентна силі, що
створюється стороннім електричним полем
\begin{equation*}
	\Efield^* = \frac{\vect{F}^*}{q} = \frac{1}{c} \vect{v} \times \Bfield
\end{equation*}
і~викликає в~перемичці довжиною \(\ell\) ЕРС індукції, рівну
\begin{equation*}
	\emf_{\text{інд}} = \oint \Efield^* \cdot d\vect{\ell} = - E_{\text{стор}} \ell = -
	\frac{1}{c} B \ell v.
\end{equation*}
Знак «\(-\)» тут пов'язаний з~тим, що \(\mathcal{E}_{\text{інд}}\) створює струм у
від'ємному напрямку (додатний напрямок задається вектором \(\Bfield\)).

%\placeholderfig{Зліва --- рамка зі струмом у магнітному полі. Перемичка може вільно ковзати, безперервно замикаючи верхній і нижній проводи рамки; справа --- напрямок сили Лоренца, що діє на рухомі позитивні заряди}{fig:5.2.1}

Оскільки \(\ell v = dS/dt\), де \(S\)~--- площа контуру, то
\begin{equation}\label{eq:faraday_law_for_moving}
	\emf_{\text{інд}} = - \frac{1}{c} \frac{d\Phi}{dt},
\end{equation}
де \(\Phi = BS\)~--- магнітний потік через контур.

Якщо є складова магнітного поля \(\Bfield_{\parallel}\), паралельна площині контуру, то
вона не призводить до появи індукційного струму вздовж контуру, оскільки відповідна
складова сили Лоренца \(\vect{F}^{(1)} = \frac{1}{c} \vect{v} \times \Bfield_{\parallel}\)
перпендикулярна до площини контуру, тобто напрямлена перпендикулярно до провідника. Це
означає, що у~виразі для потоку \(\Phi\) потрібно врахувати тільки складову
\(\Bfield_{\perp}\), тобто покласти \(\Phi = B_{\perp} S = \Bfield\cdot\vect{S}\).

Для контуру довільної форми, що рухається й~деформується в~стаціонарному, але
неоднорідному полі, ЕРС руху знаходимо, підсумовуючи по контуру сторонню силу $[\vect
		v\times\Bfield]/c$, де $\vect v$~--- швидкість елемента $d\vect\ell$:
\begin{equation*}
	\emf_\text{інд} = \frac1c\oint [\vect v\times\Bfield]\cdot d\vect\ell.
\end{equation*}
За час $dt$ елемент $d\vect\ell$ замітає площу $d\vect S = [\vect v\,dt\times
			d\vect\ell]$, тому зміна потоку дорівнює $d\Phi = \oint\Bfield\cdot[\vect v\times
		d\vect\ell]\,dt = -dt\oint[\vect v\times\Bfield]\cdot d\vect\ell$, і~знову
$\emf_\text{інд} = -\dfrac1c\dfrac{d\Phi}{dt}$, де магнітний потік визначається
інтегралом~\eqref{eq:magnetic_flux} по площі, що обмежується контуром.


\begin{Attention}
	Отже, у~розглянутому випадку \emphx{безпосередньою фізичною причиною виникнення ЕРС
		індукції є сила Лоренца}, що діє на вільні носії заряду в~рухомій перемичці. Саме вона
	«штовхає» електрони вздовж провідника, створюючи індукційний струм.

	Однак тут виникає природне запитання: що буде, якщо перейти в~систему відліку, пов'язану з
	самою перемичкою? Відносно неї провідник нерухомий, і~магнітна складова сили Лоренца, яка
	залежить від швидкості, дорівнює нулю. Якщо причина зникла, то чому тоді в~контурі все
	одно тече струм?

	Відповідь вимагає перегляду наших уявлень про електричне та магнітне поля: те, що ми
	називаємо «причиною» ЕРС, залежить від вибору системи відліку. У~системі, де перемичка
	нерухома, на її заряди діє вже \emphx{електричне} поле $\Efield' = [\vect
			v\times\Bfield]/c$ (для $v\ll c$), яке з'являється при переході в~рухому систему відліку.
	Так само й~у~досліді~1: коли магніт рухається відносно нерухомого контуру, то в~системі
	контуру магнітне поле змінюється, і~причиною є \emphx{вихрове електричне поле}, яке ми
	розглянемо в~наступному параграфі. Чому обидва описи дають однаковий результат~--- питання
	принципу відносності (гл.~\ref{Relativity}): електричне й~магнітне поля не є незалежними,
	а~їхній поділ залежить від системи відліку.
\end{Attention}

%% --------------------------------------------------------
\subsection*{Принцип роботи генераторів змінного струму}
%% --------------------------------------------------------

Явище електромагнітної індукції в~рухомих провідниках лежить в~основі роботи
\emphx{генераторів електричного струму}~--- пристроїв, без яких неможливо уявити сучасну
електроенергетику: у~них механічна енергія обертання перетворюється на електричну.

\begin{figure}[h!]\centering
	\includegraphics[width=0.45\linewidth]{AC-Generator}
	\caption{Схема найпростішого генератора змінного струму\label{pic:ACGenerator}}
\end{figure}

Найпростіший генератор складається з~рамки, кінці якої закріплені на кільцях, що
обертаються разом з~нею; до кілець щільно прилягають нерухомі щітки-контакти
(рис.~\ref{pic:ACGenerator}). Якщо рамку площею $S$ обертати в~однорідному полі $\Bfield$
з~кутовою швидкістю $\omega$, то потік крізь неї
\begin{equation*}
	\Phi(t) = BS\cos(\omega t), \qquad \frac{d\Phi}{dt} = -BS\omega\sin(\omega t),
\end{equation*}
тому в~рамці виникає ЕРС індукції
\begin{equation}\label{eq:GeneratorEMF}
	\emf_\text{інд} = \frac{BS\omega}{c}\sin(\omega t),
\end{equation}
а~якщо на щітки підключити навантаження опором $R$, у~колі потече \emphx{змінний}
індукційний струм
\begin{equation}\label{eq:GeneratorCurrent}
	I(t) = \frac{\emf_\text{інд}}{R} = \frac{BS\omega}{cR}\sin(\omega t),
\end{equation}
що змінюється за законом синуса з~кутовою частотою $\omega$, яка дорівнює кутовій
швидкості обертання рамки. Якщо рамка має $N$ витків, ЕРС збільшується в~$N$ разів. Саме
такий принцип лежить в~основі роботи практично всіх електростанцій, що генерують змінний
струм.

%% --------------------------------------------------------
\section{Електромагнітна індукція в~нерухомих провідниках}\hypertarget{induction-in-stationary-conductors}{}
%% --------------------------------------------------------

%% --------------------------------------------------------
\subsection*{Виникнення ЕРС внаслідок зміни магнітного поля}
%% --------------------------------------------------------

Розглянемо тепер принципово інший випадок: контур нерухомий і~незмінний за формою, а
змінюється саме магнітне поле, що його пронизує (рис.~\ref{pic:FieldChange}):
\begin{equation}\label{eq:FaradayFieldChange}
	\emf_\text{інд} = -\frac1c\frac{d\Phi}{dt} =
	-\frac1c\iint\limits_S\parttime{\Bfield}\cdot d\vect S.
\end{equation}

\begin{figure}[h!]\centering
	\includegraphics[width=0.5\linewidth]{Field_change}
	\caption{ЕРС індукції в~нерухомому контурі за рахунок зміни поля
		$\Bfield$\label{pic:FieldChange}}
\end{figure}

Тут виникає принципове запитання: \emphx{яка причина виникнення ЕРС}? Носії заряду в
нерухомому провіднику мають нульову швидкість, а~магнітне поле \emphx{не діє на нерухомі
	заряди} взагалі! То внаслідок чого ж вони прискорюються? Відповідь на це запитання ---
\emphx{вихрове електричне поле}~--- ми дамо трохи згодом; спершу розглянемо кілька
важливих якісних наслідків закону індукції.

%% --------------------------------------------------------
\subsection*{Правило Ленца}
%% --------------------------------------------------------

\emphx{Правило Ленца} визначає \emphx{напрям} індукційного струму: індукційний струм
завжди має такий напрям, щоб своїм власним магнітним полем \emphx{протидіяти} зміні
магнітного потоку, яка його породжує, тобто щоб послабити причину, яка збуджує цей струм
(рис.~\ref{pic:LenzRule}).

\begin{figure}[h!]\centering
	\localinput{lenz-rule.tikz}
	\caption{Правило Ленца: напрям індукційного струму залежить від знака
		$\partial\Bfield/\partial t$\label{pic:LenzRule}}
\end{figure}

\begin{Attention}
	Правило Ленца~--- не окремий незалежний закон природи, а~прямий наслідок знака
	\enquote{мінус} у~законі індукції~\eqref{eq:FaradayLawInt}, який, своєю чергою,
	диктується законом збереження енергії: якби індукційний струм \emphx{підсилював} зміну
	потоку (а~не протидіяв їй), процес наростав би лавиноподібно й~необмежено, породжуючи
	енергію з~нічого. Це нагадує загальніший \emphx{принцип Ле Шательє-Брауна}: будь-який
	зовнішній вплив, що намагається вивести систему зі стійкої рівноваги, породжує в~ній
	компенсаційні процеси, спрямовані на послаблення цього впливу.
\end{Attention}

\begin{Historical}
	Правило напряму індукційного струму сформулював у~1833--1834~рр. Е.~Х.~Ленц, а~закон
	індукції в~кількісному вигляді, через зміну магнітного потоку, записав Ф.~Нейман
	(1845~р.). Сам Фарадей користувався картиною \enquote{ліній сили}, що перетинаються
	провідником. Незалежно від Фарадея індукцію спостерігав Дж.~Генрі (1832~р.), а
	Г.~Гельмгольц (1847~р.) показав, що знак \enquote{мінус} у~законі індукції випливає із
	закону збереження енергії.
\end{Historical}

%% --------------------------------------------------------
\subsection*{Струми Фуко}
%% --------------------------------------------------------

\emphx{Струмами Фуко} (вихровими струмами) називають індукційні струми, що виникають не в
тонких лінійних провідниках, а~в~товщі суцільних масивних провідників при зміні магнітного
потоку крізь них (рис.~\ref{pic:FukoCurrents}). Вони підпорядковані тому самому правилу
Ленца: їхнє магнітне поле напрямлене так, щоб протидіяти зміні потоку, яка їх породжує.

\begin{SCfigure}[1][h!]\centering
	\includegraphics[width=0.45\linewidth]{Eddy_currents}
	\caption{Індукування вихрових струмів Фуко в~провідному диску, що обертається в~полі
		магніту\label{pic:FukoCurrents}}
\end{SCfigure}

\begin{Attention}
	Оскільки струми Фуко течуть у~товщі провідника (а~не по тонкому дроту з~відомим
	опором), вони можуть спричиняти \emphx{значне} тепловиділення (закон Джоуля-Ленца,
	гл.~\ref{SteadyCurrent})~--- саме на цьому явищі, зокрема, ґрунтується
	\emphx{індукційне нагрівання} металів. З~іншого боку, ті самі струми, взаємодіючи з
	магнітним полем (сила Ампера), створюють силу, що \emphx{гальмує} рух провідника
	відносно джерела поля~--- явище \emphx{магнітного демпфування}, яке використовують,
	наприклад, в~електромагнітних гальмах і~в~заспокоювачах коливань вимірювальних
	приладів. Швидкозмінні струми Фуко витісняються до поверхні провідника (скін-ефект,
	гл.~\ref{MaxwellEqns}), тому масивні осердя трансформаторів набирають з~ізольованих
	пластин або виготовляють із фериту, щоб зменшити втрати на нагрівання.
\end{Attention}

%% --------------------------------------------------------
\subsection*{Вихрове електричне поле}
%% --------------------------------------------------------

Повернімося до питання, яке ми залишили відкритим: що саме прискорює нерухомі заряди при
зміні $\Bfield$? Оскільки $\Phi = \iint_S \Bfield\cdot d\vect S$, а~ЕРС, за
означенням~\eqref{eq:EMFDef} (гл.~\ref{SteadyCurrent}), дорівнює $\emf = \oint_L
	\Efield\cdot d\vect\ell$, із закону індукції безпосередньо випливає:
\begin{equation}\label{eq:CurlEInt}
	\tcbhighmath{ \oint\limits_L \Efield\cdot d\vect\ell = -\frac1c\iint\limits_S
		\frac{\partial\Bfield}{\partial t}\cdot d\vect S. }
\end{equation}
Це~--- \emphx{теорема про циркуляцію} електричного поля: циркуляція напруженості
електричного поля по довільному контуру дорівнює, взятому з~протилежним знаком і
поділеному на $c$, потоку вектора $\partial\Bfield/\partial t$ крізь довільну поверхню,
натягнуту на цей контур. За теоремою Стокса, цю теорему можна записати й~у~диференціальній
формі:
\begin{equation}\label{eq:CurlEDiff}
	\tcbhighmath{ \rot\Efield = -\frac1c\frac{\partial\Bfield}{\partial t}. }
\end{equation}

Зауважимо, що теорему~\eqref{eq:CurlEInt} записано для \emphx{нерухомого} контуру. Якщо
контур рухається, до $\Efield$ у~формулі для ЕРС слід додати силу Лоренца, що припадає на
одиницю заряду: $\emf_\text{інд} = \oint\bigl(\Efield + [\vect
		v\times\Bfield]/c\bigr)\cdot d\vect\ell = -\dfrac1c\dfrac{d\Phi}{dt}$. Це~--- загальний
закон індукції~\eqref{eq:FaradayLawInt}, який охоплює обидва механізми: сила Лоренца
відповідає за ЕРС руху (попередній параграф), а~вихрове електричне поле~--- за ЕРС у
нерухомому контурі.

За Максвеллом, явище електромагнітної індукції полягає саме в~тому, що \emphx{будь-яке
	змінне магнітне поле збуджує в~просторі електричне поле}; провідники для цього взагалі не
потрібні~--- вони лише дають змогу побачити це поле через індукційний струм, який воно
викликає в~них (рис.~\ref{pic:VortexField}).

\begin{SCfigure}[1.5][h!]\centering
	\localinput{vortex-field.tikz}
	\caption{Вихрове електричне поле, породжене зростаючим у~часі магнітним
		полем\label{pic:VortexField}}
\end{SCfigure}

\begin{Attention}
	На відміну від електростатики, де завжди $\rot\Efield=0$ (поле потенціальне), у
	присутності змінного магнітного поля $\rot\Efield\neq0$. Це принципово важлива
	відмінність: індуковане електричне поле не є потенціальним~--- воно \emphx{вихрове}:
	його силові лінії ніде не починаються й~не закінчуються (на відміну від ліній
	електростатичного поля, які завжди починаються й~закінчуються на зарядах), а~замкнені
	самі на себе, подібно до силових ліній магнітного поля. Електричне поле, породжене
	зміною магнітного поля, має \emphx{зовсім іншу структуру}, ніж електростатичне, і~не
	пов'язане безпосередньо з~жодними електричними зарядами.

	Чому ж тоді це поле взагалі називають \emphx{електричним}, якщо воно має інше
	походження й~іншу конфігурацію ліній, ніж звичне статичне поле? Відповідь проста й
	фундаментальна: \emphx{вихрове поле діє на заряд точнісінько так само, як і
		електростатичне},~--- а~саме це ми з~самого початку курсу вважали \emphx{головною
		визначальною властивістю} електричного поля. Саме на прискоренні заряджених частинок
	вихровим електричним полем ґрунтується робота \emphx{бетатрона}~--- прискорювача
	електронів, у~якому змінне поле потужного електромагніта породжує вихори електричного
	поля, здатні розганяти електрони до швидкостей, близьких до швидкості світла, причому
	розгін відбувається безпосередньо у~вакуумній камері, без жодних металевих
	провідників.
\end{Attention}

%% --------------------------------------------------------
\section{Явище самоіндукції}\hypertarget{self-induction}{}
%% --------------------------------------------------------

%% --------------------------------------------------------
\subsection*{Індуктивність контуру}
%% --------------------------------------------------------

Зміна струму в~довільному контурі змінює й~створюване ним магнітне поле, а~отже~--- і
власний магнітний потік крізь цей самий контур; за законом індукції це, своєю чергою,
породжує в~контурі ЕРС. Це явище називають \emphx{самоіндукцією}
(рис.~\ref{pic:SelfInduction}).

\begin{SCfigure}[1.25][h!]\centering
	\localinput{self-induction-field.tikz}
	\caption{Контур зі струмом $I$ і~створюване ним власне магнітне
		поле\label{pic:SelfInduction}}
\end{SCfigure}

Якщо поблизу контуру немає феромагнетиків, поле $\Bfield$, а~отже, й~повний магнітний
потік $\Phi$ крізь контур пропорційні силі струму $I$ в~ньому:
\begin{equation}\label{eq:SelfInductionFlux}
	\tcbhighmath{ \Phi = \frac1c LI, }
\end{equation}
де коефіцієнт пропорційності $L$ називають \emphx{індуктивністю контуру}~--- суто
геометричною характеристикою (залежною від форми, розмірів контуру й~від магнітних
властивостей середовища навколо нього), яка не залежить від сили струму. При зміні струму
в~контурі, за законом Фарадея, виникає \emphx{ЕРС самоіндукції}:
\begin{equation}\label{eq:SelfInductionEMF}
	\tcbhighmath{ \emf_\text{сі} = -\frac{L}{c^2}\frac{dI}{dt}. }
\end{equation}

\begin{Attention}
	Знак \enquote{мінус} тут~--- знову прояв правила Ленца: ЕРС самоіндукції завжди
	напрямлена так, щоб протидіяти зміні сили струму в~контурі~--- вона гальмує зростання
	струму й~\emphx{підтримує} його спадання. Через цю властивість контур з~великою
	індуктивністю поводиться подібно до тіла з~великою механічною інерцією (масою): він
	\enquote{опирається} будь-якій швидкій зміні струму.
\end{Attention}

%% --------------------------------------------------------
\begin{Example}{Приклади обчислення індуктивності}

	\emphx{Соленоїд.} Магнітне поле всередині довгого соленоїда з~магнетиком проникності
	$\mu$, густиною намотки $n=N/l$ витків на одиницю довжини (гл.~\ref{MagFieldMedia}),
	дорівнює $B = 4\pi\mu NI/(cl)$. Повний магнітний потік крізь усі $N$ витків соленоїда
	$\Phi = NBS = \dfrac{4\pi\mu N^2}{cl}IS$; порівнюючи з
	означенням~\eqref{eq:SelfInductionFlux}, отримуємо індуктивність соленоїда:
	\begin{equation}\label{eq:SolenoidInductance}
		L = \frac{4\pi\mu N^2 S}{l}.
	\end{equation}

	\emphx{Тороїд.} Для тороїдальної котушки з~$N$ витками, прямокутним перерізом ширини $b-a$
	(де $a$ і~$b$~--- внутрішній і~зовнішній радіуси) і~висотою $d$ (див. рис.) поле за
	теоремою про циркуляцію дорівнює $B(r)=2\mu NI/(cr)$, а~потік крізь один виток ---
	$\int_a^b B\,d\,dr$; тому індуктивність дорівнює

	\begin{center}%[h!]\centering
		\includegraphics[width=0.3\linewidth]{toroid}
		%	\caption{Тороїдальна котушка\label{pic:Toroid}}
	\end{center}
	\begin{equation}\label{eq:ToroidInductance}
		L = 2\mu N^2 d\ln\left(\frac{b}{a}\right).
	\end{equation}
\end{Example}

\begin{Attention}
	В~гауссовій системі одиниць індуктивність, як неважко переконатись з
	формул~\eqref{eq:SolenoidInductance}--\eqref{eq:ToroidInductance}, вимірюється просто
	в~\emphx{сантиметрах}, $[L]=\text{см}$~--- ще одне яскраве підтвердження того, що
	індуктивність, як і~ємність (гл.~\ref{ElEnergy}),~--- суто \emphx{геометрична}
	характеристика контуру (при сталому $\mu$). Для переходу до практичних одиниць: $1\
		\text{Гн}=10^{9}\ \text{см}$.
\end{Attention}

%% --------------------------------------------------------
\subsection*{Перехідні процеси в~колі з~індуктивністю}
%% --------------------------------------------------------

\emphx{Встановлення струму в~$LR$-колі.} Розглянемо коло з~послідовно з'єднаних джерела
ЕРС $\emf$, котушки індуктивності $L$, резистора $R$ і~ключа $S$
(рис.~\ref{pic:LRTurnOn}). Після замикання ключа закон Ома для кола, $\emf +
	\emf_\text{сі} = IR$, з~урахуванням формули~\eqref{eq:SelfInductionEMF}, набуває вигляду
\begin{equation}\label{eq:LRTurnOnEq}
	\frac{L}{c^2}\frac{dI}{dt} + IR = \emf.
\end{equation}

\begin{figure}[h!]\centering
	\localinput{lr-turnon.tikz}
	\caption{Встановлення струму в~колі з~індуктивністю після замикання
		ключа\label{pic:LRTurnOn}}
\end{figure}

Це рівняння~--- повний аналог рівняння зарядки конденсатора (гл.~\ref{SteadyCurrent}), і
його розв'язок (з~початковою умовою $I(0)=0$)~--- струм, що асимптотично наближається до
стаціонарного значення $\emf/R$:
\begin{equation}\label{eq:LRTurnOnSolution}
	I(t) = \frac{\emf}{R}\left(1-e^{-t/\tau}\right), \qquad \tau = \frac{L}{c^2R}
\end{equation}
--- сталу часу $\tau$ називають \emphx{часом релаксації} $LR$-кола.

\emphx{Екстраструми при розмиканні.} Нехай тепер котушка з~індуктивністю $L$ і~власним
(внутрішнім) опором обмотки $r$ спочатку під'єднана до джерела ЕРС $\emf$ через ключ $S$ і
резистор $R$, під'єднаний паралельно котушці (рис.~\ref{pic:LRTurnOff}); у~сталому режимі
через котушку тече струм $I_0 = \emf/r$.

\begin{figure}[h!]\centering
	\localinput{lr-turnoff.tikz}
	\caption{Екстраструм розмикання в~колі з~індуктивністю\label{pic:LRTurnOff}}
\end{figure}

Після розмикання ключа $S$ (відключення джерела) магнітне поле котушки, а~отже, і~потік
крізь неї, починають спадати. Це, за законом Ленца, збуджує ЕРС самоіндукції
$\emf_\text{сі}$ і~підтримуваний нею \emphx{екстраструм розмикання}~--- індукційний струм,
що тепер тече вже через резистор $R$ (замкнене коло $L$--$r$--$R$ без джерела). Закон Ома
для цього контуру: $I(R+r) = -\dfrac{L}{c^2}\dfrac{dI}{dt}$, звідки
\begin{equation}\label{eq:LRTurnOffSolution}
	I(t) = I_0\, e^{-t/\tau'}, \qquad \tau' = \frac{L}{c^2(R+r)}.
\end{equation}

\begin{Attention}
	Якщо $R\gg r$ (типова ситуація, коли навантаження має значно більший опір, ніж провід
	самої котушки), то в~момент розмикання ЕРС самоіндукції досягає значення
	$\emf_\text{сі} = \dfrac{R+r}{r}\,\emf\, e^{-t/\tau'}\approx\dfrac{R}{r}\,\emf\,
		e^{-t/\tau'}$, яке при $R\gg r$ може \emphx{значно перевищувати} ЕРС самого джерела
	$\emf$! Саме тому раптове розмикання кола з~великою індуктивністю (наприклад,
	вимкнення потужного електромагніта чи реле) часто супроводжується електричним пробоєм
	або яскравою іскрою на контактах ключа~--- індуктивність \enquote{чинить опір} різкому
	зникненню струму настільки завзято, що готова створити на короткий час дуже високу
	напругу.
\end{Attention}

%% --------------------------------------------------------
\section{Індукція в~надпровідниках. Ефект Мейснера}\hypertarget{superconductors-induction}{}
%% --------------------------------------------------------

В~розділі~\ref{CurrentsInMedia} ми вже розглядали явище надпровідності. Для електромагнітної
індукції існування надпровідного стану цікаві: стала часу кола $\tau=L/(c^2R)$ при $R\to0$
стає нескінченною, а~індукційні струми~--- незгасаючими.

%% --------------------------------------------------------
\subsection*{Замороження потоку}
%% --------------------------------------------------------

Розглянемо замкнене кільце з~індуктивністю $L$ й~опором $R$ у~зовнішньому полі
(рис.~\ref{tikz:flux-conservation}).

Сумарний потік крізь кільце складається із зовнішнього $\Phi_\text{зовн}$ і~власного
$LI/c$, а~повна ЕРС індукції дорівнює $IR$:
\begin{equation}\label{eq:SCFluxLaw}
	-\frac1c\frac{d}{dt}\left(\Phi_\text{зовн}+\frac{LI}{c}\right)=IR.
\end{equation}
Якщо $R=0$ (надпровідне кільце), то
\begin{equation}\label{eq:SCFluxFrozen}
	\Phi_\text{зовн}+\frac{LI}{c}=\const
\end{equation}
--- повний магнітний потік крізь кільце \emphx{зберігається} (замороження потоку).
Будь-яку зміну зовнішнього потоку кільце компенсує струмом $\Delta
	I=-c\,\Delta\Phi_\text{зовн}/L$ (це правило Ленца в~найчистішому вигляді), а~якщо зовнішнє
поле зникає, у~кільці лишається \emphx{незгасаючий струм} $I=c\Phi/L$, що підтримує
колишній потік.

%---------------------------------------------------------
\begin{SCfigure}[1.5][h!]\centering
	\localinput{flux-conservation.tikz}
	\caption{Замороження магнітного потоку в~надпровідному кільці при зміні його положення
		в~зовнішньому полі: кільце підтримує колишній потік незгасаючим
		струмом\label{tikz:flux-conservation}}
\end{SCfigure}
%---------------------------------------------------------

У~дослідах зі свинцевими кільцями спаду такого струму не вдавалося виявити протягом року,
звідки випливає оцінка часу спаду понад $10^{5}$ років; для порівняння, стала часу
звичайного $LR$-кола~--- частки секунди.



%% --------------------------------------------------------
\subsection*{Ідеальний провідник чи надпровідник?}
%% --------------------------------------------------------

У~масивному зразку з~нескінченною провідністю закон Ома $\vect j=\lambda\Efield$ дає
$\Efield=0$ (інакше струм був би нескінченним), а~тоді з
$\rot\Efield=-\frac1c\,\partial\Bfield/\partial t$~\eqref{eq:CurlEDiff} маємо
$\partial\Bfield/\partial t=0$: індукція всередині ідеального провідника \emphx{не може
	змінюватися}, і~зразок зберігає те поле, що було в~ньому в~момент переходу.

Проте в~1933~р. В.~Мейснер і~Р.~Оксенфельд виявили, що для справжніх надпровідників це не
так: якщо зразок, що перебуває в~слабкому полі, охолодити нижче $T_c$, поле з~нього
\emphx{виштовхується}, і~в~товщі зразка $\Bfield=0$ незалежно від того, чи поле було
прикладене до переходу, чи після нього (\emphx{ефект Мейснера}). Це не можна пояснити
самою лише індукцією при $R\to0$, яка, навпаки, заморожувала б поле. Отже, надпровідність
--- не просто $R=0$, а~особливий стан речовини з~$\Bfield=0$ всередині.

%% --------------------------------------------------------
\subsection*{Надпровідник як ідеальний діамагнетик}
%% --------------------------------------------------------

У~надпровіднику $\Bfield=\Hfield+4\pi\vect M=0$, тому намагніченість $\vect
	M=-\Hfield/4\pi$, тобто $\chi=-1/4\pi$ і~$\mu=0$ (гл.~\ref{MagFieldMedia}). Це
найсильніший можливий діамагнетизм: у~звичайних діамагнетиків
$|\chi|\sim10^{-6}$--$10^{-5}$. Поле виштовхують поверхневі екрануючі струми, які течуть у
тонкому шарі, а~не строго по поверхні.

Товщину цього шару можна оцінити з~рівнянь Ф.~і~Г.~Лондонів (1935~р.). Для носіїв заряду
$q$ і~маси $m$ з~концентрацією $n_s$ без зіткнень $m\,d\vect v/dt=q\Efield$, тобто
$\partial\vect j/\partial t=(n_sq^2/m)\Efield$. Узявши ротор і
скориставшись~\eqref{eq:CurlEDiff}, знаходимо
\begin{equation*}
	\frac{\partial}{\partial t}\left[\rot\vect j+\frac{n_sq^2}{mc}\Bfield\right]=0.
\end{equation*}
Для ідеального провідника вираз у~дужках~--- стала величина (зберігається початкове поле).
Лондони постулювали, що в~надпровіднику він дорівнює \emphx{нулю}:
\begin{equation}\label{eq:LondonSecond}
	\rot\vect j=-\frac{n_sq^2}{mc}\Bfield.
\end{equation}
Разом з~$\rot\Bfield=4\pi\vect j/c$ (гл.~\ref{MagFieldVacuum}) і~$\divg\Bfield=0$ це дає
$\nabla^2\Bfield=\Bfield/\delta_L^2$, де
\begin{equation}\label{eq:LondonDepth}
	\tcbhighmath{ \delta_L=\sqrt{\frac{mc^2}{4\pi n_sq^2}}. }
\end{equation}
Поблизу плоскої поверхні розв'язок має вигляд $B=B_0e^{-x/\delta_L}$: поле проникає лише
на \emphx{лондонівську глибину} $\delta_L$. При $n_s\sim10^{22}$~см$^{-3}$, $q=e$, $m=m_e$
маємо $\delta_L\sim5\cdot10^{-6}$~см, тобто десятки нанометрів, що добре узгоджується з
дослідами. (Для куперівських пар $q=2e$, $m\to2m_e$, $n_s\to n/2$, а~комбінація $n_sq^2/m$
та сама.)

Розглянемо тепер надпровідну кулю радіуса $a$, вміщену в~зовнішнє однорідне магнітне поле
$\Bfield_0$ (рис.~\ref{pic:BSupercond}). Припустимо, що $a\gg\delta_L$, тобто
лондонівський шар проникнення дуже тонкий порівняно з~розмірами кулі. Тоді в~усьому об'ємі
надпровідника, за винятком тонкого приповерхневого шару, можна вважати $\Bfield=0$.

%---------------------------------------------------------
\begin{figure}[h!]\centering
	\localinput{BSupercond.tikz}
	\caption{Надпровідна куля в~зовнішньому магнітному полі: силові лінії обтікають кулю,
		не проникаючи всередину~--- поле виштовхується поверхневими екрануючими
		струмами\label{pic:BSupercond}}
\end{figure}
%---------------------------------------------------------

Між поведінкою провідника в~електростатичному полі та надпровідника в~магнітному полі
існує глибока аналогія. В~обох випадках матеріал \emphx{екранує} зовнішнє поле: провідник
перерозподіляє вільні заряди так, щоб усередині нього зникло електричне поле,

\begin{equation}
	\Efield_{\mathrm{in}}=0,
\end{equation}
а~надпровідник створює екрануючі струми, які виштовхують магнітне поле з~його об'єму,

\begin{equation}
	\Bfield_{\mathrm{in}}\simeq0, \qquad a\gg\delta_L.
\end{equation}

В~обох випадках екранування забезпечується перерозподілом заряджених носіїв: у~звичайному
провіднику виникає поверхнева густина заряду, а~в~надпровіднику~--- поверхневий струм у
тонкому шарі проникнення. Ці індуковані джерела створюють поле, яке компенсує зовнішнє
поле всередині матеріалу. Проте механізми екранування фізично різні: електростатичне
екранування є наслідком перерозподілу вільних зарядів до стану рівноваги, тоді як магнітне
екранування надпровідника є проявом ефекту Мейснера. Тому надпровідник не слід просто
ототожнювати з~ідеальним провідником.

Спільним є й~математичний апарат. Обидві задачі~--- про точковий заряд біля провідної
поверхні та про магнітний диполь біля надпровідної поверхні~--- розв'язують
\emphx{методом зображень}. Підставою для цього є те, що в~обох задачах поле поза
матеріалом задовольняє рівняння Лапласа, а~на поверхні матеріалу задана гранична
умова, яку можна відтворити, замінивши матеріал фіктивним джерелом за поверхнею. Для
провідника така умова~--- $\Efield_\tau=0$, і~її забезпечує фіктивний заряд
протилежного знака; для надпровідника~--- $B_n=0$, і~її забезпечує фіктивний диполь
протилежного напрямку.

%% --------------------------------------------------------
\subsection*{Левітація}
%% --------------------------------------------------------

Виштовхування поля означає, що на надпровідник діє сила. Скористаймося методом зображень:
магніт над плоскою надпровідною пластиною поводиться так, ніби на такій самій глибині під
поверхнею лежить його \enquote{дзеркальне відображення}~--- диполь, спрямований
протилежно. Два співвісні антипаралельні диполі на відстані $2h$ відштовхуються з~енергією
$2p_m^2/(2h)^3$; оскільки відображення уявне, енергія взаємодії вдвічі менша:
\begin{equation}\label{eq:LevitationForce}
	U=\frac{p_m^2}{8h^3},\qquad F=-\frac{dU}{dh}=\frac{3p_m^2}{8h^4},
\end{equation}
де $h$~--- висота диполя над поверхнею. Прирівнявши $F=mg$, знаходимо висоту левітації
$h=(3p_m^2/8mg)^{1/4}$ (задача~\ref{prob:MagnetLevitation}).

% \begin{Attention}
% 	Левітація не суперечить \emphx{теоремі Ірншоу} (1842~р.), за якою в статичному полі лише зарядів і постійних магнітів стійка
% 	рівновага неможлива. Тут енергія $U\propto H^2$, а для поля без джерел ($\rot\Hfield=0$, $\divg\Hfield=0$)
% 	$\nabla^2H^2=2\sum(\nabla H_i)^2\ge0$, тому мінімум енергії допустимий: діамагнетики й надпровідники \enquote{обходять} теорему.
% \end{Attention}

\begin{Historical}
	За легендою, залізна труна пророка Магомета висить у~повітрі без опори між магнітами. Для
	системи, що складається лише зі звичайних постійних магнітів, стійка статична левітація
	неможлива внаслідок теореми Ерншоу (див. розділ~\ref{SteadyEfield}). Надпровідник, однак,
	дозволяє реалізувати стійку магнітну левітацію: його екрануючі струми не є наперед
	заданими джерелами, а~виникають у~відповідь на зовнішнє магнітне поле та забезпечують
	ефект Мейснера. У~результаті взаємодія магніту з~цими індукованими струмами може створити
	стійке положення рівноваги. У~1945~р. В.~К.~Аркадьєв продемонстрував левітацію постійного
	магніту над надпровідною свинцевою чашею, перетворивши легенду про \enquote{труну
		Магомета} на лабораторний експеримент.
\end{Historical}

%% --------------------------------------------------------*
\subsection*{Квантування потоку}
%% --------------------------------------------------------*

У~надпровідному кільці потік не лише зберігається, а~й~\emphx{квантується}:

\begin{equation}
	\Phi=n\Phi_0, \qquad \Phi_0=\frac{\pi\hbar c}{|e|}
	\approx2{,}07\cdot10^{-7}~\mathrm{Гс}\cdot\mathrm{см}^2.
\end{equation}

Квант потоку відповідає заряду носіїв \(2e\) (куперівські пари). Його експериментально
виміряли в~1961~р. (Р.~Дівер і~У.~Фейрбенк; Б.~Долл і~М.~Небауер), а~на квантуванні потоку
ґрунтуються надчутливі магнітометри SQUID.

%%% --------------------------------------------------------
\subsection*{Критичні параметри}
%%% --------------------------------------------------------


Надпровідний стан має межі: він існує лише за певних температур, магнітних полів і
струмів. Для надпровідників першого роду при \(T\to0\) критичне поле \(H_c\) становить
сотні гаусів: для свинцю \(H_c\approx800~\mathrm{Гс}\), для алюмінію
\(H_c\approx100~\mathrm{Гс}\).

Надпровідність може бути зруйнована і~струмом, якщо його власне магнітне поле біля
поверхні досягає критичного значення. Для дроту радіуса \(r\) правило Сілсбі дає

\begin{equation}
	I_c=\frac{crH_c}{2}.
\end{equation}

Надпровідники другого роду поводяться інакше: між нижнім і~верхнім критичними полями
\(H_{c1}\) та \(H_{c2}\) магнітне поле проникає в~надпровідник у~вигляді
\emphx{квантованих вихорів}. Закріплення вихорів на дефектах кристалічної структури
перешкоджає їхньому руху і~тому відіграє важливу роль у~стійкості левітації.

Сплави \ce{NbTi} та \ce{Nb3Sn} здатні витримувати поля понад \(20~\mathrm{Тл}\), тому з
них виготовляють обмотки сильних магнітів для магнітно-резонансних томографів і
прискорювачів. Високотемпературні керамічні надпровідники, зокрема YBCO з
\(T_c\approx93~\mathrm{К}\), дозволяють реалізовувати надпровідний стан при охолодженні
рідким азотом.



%% --------------------------------------------------------
\section*{Підсумки}
\phantomsection \addcontentsline{toc}{section}{Підсумки}
%% --------------------------------------------------------

Явище електромагнітної індукції~--- виникнення ЕРС у~контурі при зміні охоплюваного ним
магнітного потоку~--- остаточно поєднує електрику та магнетизм у~єдину теорію. Кількісно
воно описується законом $\emf_\text{інд}=-\dfrac1c\dfrac{d\Phi}{dt}$, а~знак
\enquote{мінус} (правило Ленца) є наслідком закону збереження енергії: індукційний струм
завжди протидіє зміні потоку, що його породжує.

Потік може змінюватися внаслідок \emphx{руху} контуру (зміна площі чи орієнтації) або
внаслідок зміни \emphx{поля}. У~першому випадку ЕРС створює сила Лоренца, що діє на носії
заряду в~провіднику, що рухається (ЕРС руху); на цьому ґрунтується робота генераторів. У
другому випадку носії заряду нерухомі, і~їх прискорює \emphx{вихрове електричне поле}
$\rot\Efield=-\dfrac1c\dfrac{\partial\Bfield}{\partial t}$: воно не є потенціальним, але
діє на заряд так само як і~електростатичне (бетатрон, індукційне нагрівання, струми Фуко).
Поділ на ці два механізми залежить від системи відліку, а~загальний закон індукції охоплює
обидва: $\emf_\text{інд}=\oint(\Efield+[\vect v\times\Bfield]/c)\cdot
	d\vect\ell=-\dfrac1c\dfrac{d\Phi}{dt}$. Сила Лоренца (гл.~\ref{MagFieldVacuum}) демонструє
фундаментальний зв'язок магнітного поля з~рухом зарядів, а~вихрове поле~--- з~тим, що
змінне магнітне поле саме породжує електричне.

Змінний струм у~контурі створює змінний власний потік, і~тому в~ньому виникає ЕРС
самоіндукції $\emf_\text{сі}=-\dfrac{L}{c^2}\dfrac{dI}{dt}$. Індуктивність $L$ залежить
від форми контуру й~середовища, а~не від сили струму; вона відіграє роль \enquote{інерції}
кола. Унаслідок цього струм у~$LR$-колі встановлюється за експонентою з~часом релаксації
$\tau=L/(c^2R)$, а~при розмиканні кола з~великою індуктивністю виникають екстраструми та
ЕРС, що можуть багаторазово перевищувати ЕРС джерела.

У~надпровідників ($R=0$) сумарний потік крізь замкнене кільце зберігається,
$\Phi_\text{зовн}+LI/c=\const$, а~індукційні струми не згасають. Ефект Мейснера показує,
що надпровідник~--- не просто ідеальний провідник: у~товщі зразка $\Bfield=0$ за будь-якої
передісторії, тобто це ідеальний діамагнетик ($\chi=-1/4\pi$, $\mu=0$); поле проникає лише
на лондонівську глибину $\delta_L$. Виштовхування поля породжує силу, що робить можливою
левітацію магніту над надпровідником.

% Основні формули розділу:
% \begin{itemize}
% 	\item закон індукції $\emf_\text{інд}=-\dfrac1c\dfrac{d\Phi}{dt}$, $\Phi=\iint\Bfield\cdot d\vect S$; ЕРС руху $\emf_\text{інд}=\dfrac1c\oint[\vect v\times\Bfield]\cdot d\vect\ell$, для перемички $\emf_\text{інд}=-B\ell v/c$;
% 	\item генератор: $\emf_\text{інд}=\dfrac{NBS\omega}{c}\sin\omega t$;
% 	\item теорема про циркуляцію $\oint\Efield\cdot d\vect\ell=-\dfrac1c\iint\dfrac{\partial\Bfield}{\partial t}\cdot d\vect S$, $\rot\Efield=-\dfrac1c\dfrac{\partial\Bfield}{\partial t}$;
% 	\item самоіндукція: $\Phi=LI/c$, $\emf_\text{сі}=-\dfrac{L}{c^2}\dfrac{dI}{dt}$; соленоїд $L=4\pi\mu N^2S/l$, тороїд $L=2\mu N^2d\ln(b/a)$;
% 	\item надпровідник: $\Phi_\text{зовн}+LI/c=\const$, $\Bfield=0$ у товщі, $\chi=-1/4\pi$, $\delta_L=\sqrt{mc^2/(4\pi n_sq^2)}$;
% 	\item $LR$-коло: $I=\dfrac{\emf}{R}(1-e^{-t/\tau})$, $\tau=\dfrac{L}{c^2R}$; розмикання: $I=I_0e^{-t/\tau'}$, $\tau'=\dfrac{L}{c^2(R+r)}$.
% \end{itemize}
% Наступний логічний крок --- дослідити, як робота проти індукованого електричного поля при встановленні струму акумулюється в просторі, тобто розрахувати \emphx{енергію магнітного поля}.

%% --------------------------------------------------------
\section*{Задачі}
\phantomsection \addcontentsline{toc}{section}{Задачі}
%% --------------------------------------------------------

\begin{problem}\label{prob:RodOnRails}
	\emphx{Перемичка на рейках.} Перемичка завдовжки $\ell=20$~см рівномірно ковзає по
	рейках зі швидкістю $v=5$~м/с в~однорідному полі $B=500$~Гс, перпендикулярному до
	площини рейок. Опір замкненого контуру $R=0{,}1$~Ом. Знайдіть ЕРС індукції, струм,
	силу Ампера й~потужність, яку потрібно витрачати на рух перемички. Покажіть, що ця
	потужність дорівнює тепловій потужності $I^2R$. (Використайте $1$~СГСЕ потенціалу
	$=300$~В, $1$~А $=3\cdot10^{9}$~СГСЕ струму.) \emphx{Відповідь:} $\emf=B\ell
		v/c\approx1{,}7\cdot10^{-4}$~СГСЕ $=0{,}05$~В; $I=0{,}5$~А; $F=I\ell B/c=500$~дин
	$=5$~мН; $P=Fv=25$~мВт $=\emf I=I^2R$.
\end{problem}

\begin{problem}\label{prob:FaradayDisk}
	\emphx{Диск Фарадея.} Мідний диск радіуса $a=10$~см обертається з~частотою $50$~об/с в
	однорідному полі $B=1000$~Гс, перпендикулярному до диска. Знайдіть ЕРС між центром і
	краєм диска, якщо до цих точок дотикаються ковзні контакти. \emphx{Відповідь:} елемент
	радіуса $dr$ має швидкість $v=\omega r$, тому $\emf=\dfrac1c\int_0^a\omega
		rB\,dr=\dfrac{\omega Ba^2}{2c}\approx5{,}2\cdot10^{-4}$~СГСЕ $\approx0{,}16$~В.
\end{problem}

\begin{problem}\label{prob:GeneratorNumbers}
	\emphx{Генератор.} Рамка з~$N=100$ витків площею $S=200$~см$^2$ обертається з~частотою
	$50$~об/с в~однорідному полі $B=1000$~Гс. Знайдіть амплітуду ЕРС. Якщо рамка замкнена
	на опір $R=10$~Ом, знайдіть амплітуду струму й~середню потужність, яку потрібно
	витрачати на обертання (втрати на тертя не враховуйте). \emphx{Відповідь:}
	$\emf_m=\dfrac{NBS\omega}{c}\approx0{,}21$~СГСЕ $\approx63$~В; $I_m\approx6{,}3$~А;
	$\langle P\rangle=\dfrac{\emf_m^2}{2R}\approx0{,}2$~кВт (за законом збереження енергії
	--- рівна середній електричній потужності).
\end{problem}

\begin{problem}\label{prob:VortexCylinder}
	\emphx{Вихрове поле.} Однорідне магнітне поле $B(t)$, напрямлене вздовж осі, заповнює
	циліндричну область радіуса $a$. Знайдіть напруженість вихрового електричного поля на
	відстані $r$ від осі при $r<a$ та $r>a$. Оцініть її при $r=a=5$~см і~$\dot{B}=10^{3}$~Гс/с.
	Чому поле поза циліндром відмінне від нуля, хоча там $B=0$?
	\emphx{Відповідь:} із симетрії лінії $\Efield$~--- кола навколо осі; $2\pi
		rE=-\dfrac1c\dot B\pi r^2$ при $r<a$, тобто $E=-\dfrac{r}{2c}\dot B$; при $r>a$ потік
	дорівнює $\dot B\pi a^2$, тобто $E=-\dfrac{a^2}{2cr}\dot B$. При $r=a=5$~см
	$|E|\approx8{,}3\cdot10^{-8}$~СГСЕ $\approx2{,}5\cdot10^{-5}$~В/см. Циркуляція
	$\Efield$ по кільцю поза циліндром визначається потоком крізь \emphx{весь} циліндр,
	тому поле відмінне від нуля й~там, де $B=0$.
\end{problem}

\begin{problem}\label{prob:BetatronTurn}
	\emphx{Бетатрон.} Електрон рухається в~бетатроні по колу радіуса $R=30$~см. Потік
	магнітного поля крізь орбіту зростає зі сталою швидкістю $100$~Вб/с
	($=10^{10}$~Гс$\cdot$см$^2$/с). Яку енергію набуває електрон за один оберт? Скільки
	обертів і~який шлях потрібні, щоб набрати енергію $10$~МеВ (вважайте швидкість
	електрона близькою до $c$)? \emphx{Відповідь:}
	$\emf=\dfrac1c\dfrac{d\Phi}{dt}\approx0{,}33$~СГСЕ $=100$~В, тому за оберт електрон
	набуває $e\emf=100$~еВ; для $10$~МеВ потрібно $\approx10^{5}$ обертів, що відповідає
	шляху $\approx2\pi R\cdot10^5\approx190$~км.
\end{problem}

\begin{problem}\label{prob:FallingRod}
	\emphx{Падаюча перемичка.} Провідна перемичка масою $m=10$~г і~довжиною $\ell=20$~см
	ковзає без тертя вертикально вниз по двох вертикальних рейках у~горизонтальному
	однорідному полі $B=1000$~Гс, перпендикулярному до площини рейок. Рейки замкнено вгорі
	резистором опором $R=0{,}01$~Ом (опором іншого не враховуйте). Знайдіть усталену
	швидкість падіння та покажіть, що потужність сили тяжіння дорівнює теплу, що
	виділяється в~контурі. ($1$~Ом $\approx1{,}1\cdot10^{-12}$~с/см.) \emphx{Відповідь:}
	$mg=F_\text{А}=\dfrac{B^2\ell^2v}{c^2R}$, тому
	$v=\dfrac{mgRc^2}{B^2\ell^2}\approx2{,}5$~м/с; $mgv=I^2R$.
\end{problem}

% \begin{problem}\label{prob:SolenoidL}
% 	\emphx{Індуктивність соленоїда.} Довгий соленоїд без осердя має $N=1000$ витків, довжину $l=50$~см і радіус $r=2$~см. Знайдіть його індуктивність (у см і в генрі) та ЕРС самоіндукції при швидкості зміни струму $100$~А/с.
% 	\emphx{Відповідь:} $L=\dfrac{4\pi N^2S}{l}=\dfrac{4\pi^2N^2r^2}{l}\approx3{,}2\cdot10^{6}$~см $\approx3{,}2$~мГн; $|\emf_\text{сі}|=L\left|\dfrac{dI}{dt}\right|\approx0{,}32$~В.
% \end{problem}

\begin{problem}\label{prob:LenzFromEnergy}
	\emphx{Знак «мінус» у~законі Фарадея із закону збереження енергії.} Замкнений контур
	опором $R$ перебуває в~зовнішньому магнітному полі; потік крізь контур $\Phi(t)$
	змінюється з~часом, і~в~контурі тече індукційний струм $I$. Припустимо, що закон
	індукції має вигляд $\emf = s\,\dfrac{d\Phi}{dt}$, де $s$~--- невідомий знак ($+1$ або
	$-1$). Скориставшись законом збереження енергії, знайдіть $s$. \emphx{Вказівка:}
	порівняйте потужність, яку зовнішнє поле передає струму, із потужністю, що виділяється
	в~контурі як джоулеве тепло. Розгляньте випадок, коли потік зростає.
\end{problem}

\begin{problem}\label{prob:LRSwitchOn}
	\emphx{Встановлення струму.} Котушка з~індуктивністю $L=0{,}5$~Гн і~опором $R=50$~Ом
	підключається до джерела з~ЕРС $\emf=100$~В. Знайдіть сталу часу, стаціонарний струм,
	струм через $5$~мс після замикання та час, за який струм досягне $90\%$ стаціонарного.
	(Розрахунок можна вести в~СІ, де $\tau=L/R$.) \emphx{Відповідь:} $\tau=L/R=10$~мс;
	$I_\infty=\emf/R=2$~А; $I(5\text{~мс})=I_\infty(1-e^{-0{,}5})\approx0{,}79$~А;
	$t_{90}=\tau\ln10\approx23$~мс.
\end{problem}

\begin{problem}\label{prob:OpenCircuitEMF}
	\emphx{Екстраструм розмикання.} Котушка з~$L=1$~Гн і~опором $r=1$~Ом підключена через
	ключ до джерела з~$\emf=10$~В; паралельно котушці ввімкнено резистор $R=1$~кОм.
	Знайдіть струм у~котушці до розмикання, максимальну ЕРС самоіндукції після розмикання
	ключа та час спаду струму. Порівняйте з~$\emf$. \emphx{Відповідь:} $I_0=\emf/r=10$~А;
	$\emf_\text{сі}(0)=I_0(R+r)\approx10$~кВ, тобто у~$\sim10^{3}$ разів більше за $\emf$;
	$\tau'=L/(R+r)\approx1$~мс.
\end{problem}

\begin{problem}\label{prob:MeissnerSphere}
	\emphx{Поле навколо надпровідної кулі.} Надпровідна куля радіуса $a$ вміщена в
	однорідне зовнішнє поле $\Bfield_0$. У~наближенні ідеального діамагнетика ($\mu=0$)
	знайдіть розподіл поля зовні кулі та переконайтеся, що нормальна складова поля на її
	поверхні дорівнює нулю. \emphx{Відповідь:} $\Phi = B_0 r\cos\theta +
		\dfrac{B_0a^3}{2r^2}\cos\theta$; $B_r = B_0\cos\theta\left(1-\dfrac{a^3}{r^3}\right)$,
	$B_\theta = -B_0\sin\theta\left(1+\dfrac{a^3}{2r^3}\right)$; при $r=a$ маємо $B_r=0$,
	$B_\theta=-\dfrac32B_0\sin\theta$.
\end{problem}


\begin{problem}\label{prob:FrozenFlux}
	\emphx{Замороження потоку.} Надпровідне кільце радіуса $a=5$~см з~індуктивністю
	$L=250$~см помістили в~нормальному стані в~однорідне поле $B_0=100$~Гс,
	перпендикулярне до площини кільця, охолодили нижче $T_c$, а~потім поле вимкнули.
	Знайдіть струм у~кільці. Скільком квантам потоку $\Phi_0=\pi\hbar
		c/|e|\approx2{,}07\cdot10^{-7}$~Гс$\cdot$см$^2$ відповідає збережений потік?
	\emphx{Відповідь:} $\Phi=B_0\pi a^2\approx7{,}9\cdot10^{3}$~Гс$\cdot$см$^2$;
	за~\eqref{eq:SCFluxFrozen} $I=c\Phi/L\approx9{,}4\cdot10^{11}$~СГСЕ
	$\approx3\cdot10^{2}$~А; $\Phi/\Phi_0\approx4\cdot10^{10}$ квантів.
\end{problem}

\begin{problem}\label{prob:SilsbeeCurrent}
	\emphx{Критичний струм.} Свинцевий дріт радіуса $r=1$~мм лишається надпровідним, поки
	поле на його поверхні не досягне $H_c\approx800$~Гс (правило Сілсбі). Знайдіть
	максимальний струм без втрати надпровідності. \emphx{Відповідь:}
	$H(r)=\dfrac{2I}{cr}$, тому $I_c=\dfrac{crH_c}{2}\approx1{,}2\cdot10^{12}$~СГСЕ
	$\approx4\cdot10^{2}$~А.
\end{problem}

\begin{problem}\label{prob:MagnetLevitation}
	\emphx{Левітація магніту.} Малий магніт із магнітним моментом
	$p_m=10^{3}$~Гс$\cdot$см$^3$ (приблизно $1$~см$^3$ неодимового магніту) і~масою $8$~г,
	зорієнтований вертикально, розташований над плоскою надпровідною пластиною. На якій
	висоті він левітує? Скористайтеся~\eqref{eq:LevitationForce} і~вважайте дипольне
	наближення прийнятним. \emphx{Відповідь:}
	$h=\left(\dfrac{3p_m^2}{8mg}\right)^{1/4}\approx2{,}6$~см.
\end{problem}

\begin{problem}\label{prob:LondonDepth}
	\emphx{Глибина проникнення.} Оцініть лондонівську глибину проникнення для
	$n_s=10^{22}$~см$^{-3}$ (носії~--- електрони, $q=e$, $m=m_e$). У~скільки разів
	зменшиться поле на відстані $5\delta_L$ від поверхні надпровідника? \emphx{Відповідь:}
	$\delta_L=\sqrt{\dfrac{m_ec^2}{4\pi n_se^2}}\approx5\cdot10^{-6}$~см $=50$~нм;
	$B(5\delta_L)=B_0e^{-5}$, тобто поле зменшується приблизно у~$150$ разів.
\end{problem}

\begin{Questions}
	\begin{quizlist}
		\item Сформулюйте суть явища електромагнітної індукції. Опишіть три групи дослідів
		Фарадея (рис.~\ref{pic:FaradayExperiments}). Що в~них спільного і~чому важливою є
		лише відносна швидкість магніту та котушки?
		\item Коли в~контурі виникає індукційний струм, а~коли ні? Чому сама причина зміни
		магнітного потоку не впливає на існування струму?
		\item Запишіть закон електромагнітної індукції в~інтегральній формі. Що означає знак
		«мінус» і~якими трьома способами може змінюватися магнітний потік крізь контур?
		\item Виведіть вираз для роботи сили Ампера при зміщенні перемички на $dx$
		(рис.~\ref{tikz:frame_with_rod}) і~покажіть, що $dA = \frac{I}{c}\,d\Phi$. Яка
		складова поля $\Bfield$ виконує роботу, а~яка ні?
		\item Магнітне поле не виконує роботи, але при переміщенні перемички зі струмом робота
		виконується. Як розв'язується ця суперечність? Що таке складові швидкості
		$\vect{u}_{\perp}$ та $\vect{u}_{\parallel}$ і~яку роль відіграють сторонні ЕРС?
		\item Що є безпосередньою причиною ЕРС індукції в~рухомій перемичці
		(рис.~\ref{tikz:move_rod})? Виведіть вираз $\emf_\text{інд} = -B\ell v/c$ і
		покажіть, що він дає закон Фарадея.
		\item Чому в~рухомому провіднику причиною ЕРС є сила Лоренца, а~в~системі відліку,
		пов'язаній із самим провідником,~--- вихрове електричне поле? Про що це свідчить
		щодо співвідношення електричного й~магнітного полів?
		\item Поясніть принцип роботи найпростішого генератора змінного струму
		(рис.~\ref{pic:ACGenerator}). Виведіть формулу для ЕРС та струму в~рамці, що
		обертається в~однорідному полі.
		\item Яка причина ЕРС індукції в~нерухомому контурі, якщо магнітне поле не діє на
		нерухомі заряди?
		\item Сформулюйте правило Ленца (рис.~\ref{pic:LenzRule}). Чому воно є наслідком
		закону збереження енергії? Що сталося б, якби індукційний струм підсилював зміну
		потоку?
		\item Що таке струми Фуко? Наведіть приклади їх корисного (індукційне нагрівання,
		магнітне демпфування) та шкідливого прояву.
		\item Як записується загальний закон індукції для контуру, що рухається в~змінному
		полі? Який внесок дають сила Лоренца й~вихрове електричне поле?
		\item Запишіть теорему про циркуляцію електричного поля в~інтегральній та
		диференціальній формі. У~чому принципова відмінність вихрового електричного поля від
		електростатичного? Чому його все ж називають електричним?
		\item Що таке самоіндукція та індуктивність контуру? Від чого вона залежить, а~від
		чого ні? Запишіть ЕРС самоіндукції та поясніть її інерційну роль у~колі.
		\item Виведіть індуктивність довгого соленоїда. У~яких одиницях вимірюється
		індуктивність у~гаусовій системі і~чому це підтверджує її геометричний характер?
		\item Виведіть закон зростання струму в~$LR$-колі після замикання ключа. Що таке час
		релаксації $\tau$ і~від чого він залежить?
		\item Що таке екстраструм розмикання? Чому при $R\gg r$ ЕРС самоіндукції в~момент
		розмикання може значно перевищувати ЕРС джерела і~до яких наслідків це призводить?
		\item Що таке замороження потоку в~замкненому надпровідному контурі? Запишіть закон
		збереження потоку і~поясніть, чому індукційний струм у~ньому не згасає.
		\item Чим надпровідник відрізняється від ідеального провідника? У~чому суть ефекту
		Мейснера?
		\item Чому надпровідник називають ідеальним діамагнетиком? Чому дорівнюють $\chi$ і
		$\mu$? Що таке лондонівська глибина проникнення?
		\item Поясніть левітацію магніту над надпровідником методом зображень. Чому вона
		можлива, хоча за теоремою Ірншоу постійні магніти не можуть стійко зависати?
	\end{quizlist}
\end{Questions}